% =============================================================================
% Frontend-Backend.tex – React-Stack Grundlagen
% Einbindung: % =============================================================================
% Frontend-Backend.tex – React-Stack Grundlagen
% Einbindung: % =============================================================================
% Frontend-Backend.tex – React-Stack Grundlagen
% Einbindung: % =============================================================================
% Frontend-Backend.tex – React-Stack Grundlagen
% Einbindung: \input{content/Frontend-Backend}
% =============================================================================

\section{React-Stack Grundlagen}
\label{sec:react-stack}

% -----------------------------------------------------------------------------
\subsection{Begriffsrahmen und Zielsetzung}
\label{subsec:react-begriffsrahmen}

Du entwickelst das Selection-Panel-Frontend (Browser-UI), das über WebSocket mit
dem Python-Backend kommuniziert. Der Stack ersetzt manuelles HTML/CSS/JS durch
ein wartbares, typsicheres System.

\begin{table}[H]
\centering
\caption{Zentrale Begriffe des React-Stacks}
\label{tab:react-begriffe}
\begin{tabularx}{\textwidth}{@{}l X@{}}
\toprule
\textbf{Begriff} & \textbf{Definition} \\
\midrule
React        & JavaScript-Bibliothek für komponentenbasierte UIs \\
TypeScript   & JavaScript-Erweiterung mit statischer Typisierung \\
Vite         & Build-Tool und Entwicklungsserver (sprich: \glqq viit\grqq, französisch für \glqq schnell\grqq) \\
ESLint       & Statische Code-Analyse (findet Fehler vor Laufzeit) \\
Prettier     & Code-Formatter (einheitlicher Stil) \\
React Router & Client-seitige Navigation ohne Seitenneuladen \\
Tailwind CSS & Utility-First CSS-Framework \\
\bottomrule
\end{tabularx}
\end{table}

% -----------------------------------------------------------------------------
\subsection{React -- Komponenten statt DOM-Manipulation}
\label{subsec:react-komponenten}

Bei klassischem JavaScript manipulierst du das DOM direkt:

\begin{lstlisting}[style=base,numbers=none,caption={Klassische DOM-Manipulation},label={lst:dom-klassisch}]
// Klassisch: imperativ, fehleranfaellig bei Zustandsaenderungen
document.getElementById('led-5').classList.add('active');
document.getElementById('status').textContent = 'LED 5 an';
\end{lstlisting}

Bei 100~LEDs und komplexen Zuständen wird dieser Ansatz unwartbar. Wenn sich ein
Wert ändert, musst du manuell alle betroffenen DOM-Elemente aktualisieren.

React kehrt das Prinzip um -- du beschreibst den \emph{gewünschten Zustand},
React berechnet die nötigen DOM-Änderungen:

\begin{lstlisting}[style=base,numbers=left,caption={React-Komponente (deklarativ)},label={lst:react-komponente}]
// React: deklarativ, Zustand -> UI-Ableitung
function LedButton({ id, isActive, onPress }: LedProps) {
  return (
    <button
      className={isActive ? 'led active' : 'led'}
      onClick={() => onPress(id)}
    >
      LED {id}
    </button>
  );
}
\end{lstlisting}

\paragraph{Kernkonzept -- Komponenten}

Eine Komponente ist eine Funktion, die Props (Eingabeparameter) erhält und JSX
(HTML-ähnliche Syntax) zurückgibt. Komponenten sind wiederverwendbar und isoliert.

\paragraph{Kernkonzept -- State und Hooks}

Lokaler Zustand wird mit \code{useState} verwaltet:

\begin{lstlisting}[style=base,numbers=left,caption={State-Management mit useState},label={lst:usestate}]
function LedPanel() {
  const [activeIds, setActiveIds] = useState<Set<number>>(new Set());

  const handlePress = (id: number) => {
    setActiveIds(prev => {
      const next = new Set(prev);
      next.has(id) ? next.delete(id) : next.add(id);
      return next;
    });
  };

  return (
    <div className="grid">
      {Array.from({ length: 100 }, (_, i) => (
        <LedButton
          key={i}
          id={i}
          isActive={activeIds.has(i)}
          onPress={handlePress}
        />
      ))}
    </div>
  );
}
\end{lstlisting}

\begin{infobox}[Konsequenz für Selection Panel]
Wenn der ESP32 über WebSocket meldet \glqq LED 42 aktiv\grqq, aktualisierst du nur
den State. React rendert automatisch nur die betroffene Komponente neu -- nicht
die anderen 99~LEDs.
\end{infobox}

% -----------------------------------------------------------------------------
\subsection{TypeScript -- Fehler vor Laufzeit finden}
\label{subsec:typescript}

JavaScript-Fehler treten erst zur Laufzeit auf:

\begin{lstlisting}[style=base,numbers=none,caption={JavaScript ohne Typisierung},label={lst:js-untyped}]
// JavaScript: Fehler erst im Browser sichtbar
function setLedState(id, state) {
  websocket.send({ led: id, state: state }); // Typo? Falscher Typ?
}
setLedState("5", "an"); // Kompiliert, crasht zur Laufzeit
\end{lstlisting}

TypeScript prüft Typen beim Kompilieren:

\begin{lstlisting}[style=base,numbers=left,caption={TypeScript mit Typisierung},label={lst:ts-typed}]
// TypeScript: Fehler im Editor, bevor du speicherst
interface LedCommand {
  led: number;
  state: boolean;
}

function setLedState(id: number, state: boolean): void {
  const cmd: LedCommand = { led: id, state };
  websocket.send(JSON.stringify(cmd));
}

setLedState("5", "an"); // Fehler: string ist nicht number
setLedState(5, true);   // korrekt
\end{lstlisting}

\paragraph{Wichtige Typen für dein Projekt}

\begin{lstlisting}[style=base,numbers=left,caption={WebSocket-Nachrichten typisieren},label={lst:ws-types}]
// WebSocket-Nachrichten typisieren
type WsMessage =
  | { type: 'button_press'; id: number }
  | { type: 'led_update'; id: number; state: boolean }
  | { type: 'audio_start'; file: string };

// Props fuer Komponenten
interface LedGridProps {
  leds: Map<number, boolean>;
  onLedClick: (id: number) => void;
}
\end{lstlisting}

\begin{tipbox}[Konsequenz]
Der Compiler fängt Tippfehler, falsche Argumentreihenfolgen und fehlende
Eigenschaften ab. Bei 100~Buttons mit verschiedenen Zuständen reduziert das
die Debugging-Zeit erheblich.
\end{tipbox}

% -----------------------------------------------------------------------------
\subsection{Vite -- Build-Tool und Entwicklungsserver}
\label{subsec:vite}

Browser verstehen weder TypeScript noch JSX direkt. Der Code muss transformiert werden:

\begin{lstlisting}[style=base,numbers=none]
Quellcode (.tsx)  ->  [Transformation]  ->  Browser-Code (.js)
\end{lstlisting}

Ältere Tools wie Webpack benötigen dafür mehrere Sekunden bis Minuten.
Vite nutzt ES-Module (native Browser-Imports) und transformiert nur geänderte Dateien:

\begin{lstlisting}[style=shell,caption={Vite-Projekt erstellen},label={lst:vite-init}]
npm create vite@latest pi-react -- --template react-ts
cd pi-react
npm install
npm run dev   # Startet localhost:5173, Hot Module Replacement
\end{lstlisting}

\paragraph{Zwei Modi}

\begin{table}[H]
\centering
\caption{Vite Build-Modi}
\label{tab:vite-modi}
\begin{tabularx}{\textwidth}{@{}l l X@{}}
\toprule
\textbf{Modus} & \textbf{Befehl} & \textbf{Ergebnis} \\
\midrule
Entwicklung & \code{npm run dev}   & Lokaler Server mit Live-Reload, Source Maps \\
Produktion  & \code{npm run build} & Optimierte statische Dateien in \code{dist/} \\
\bottomrule
\end{tabularx}
\end{table}

\paragraph{Konfiguration (vite.config.ts)}

\begin{lstlisting}[style=base,numbers=left,caption={Vite-Konfiguration mit WebSocket-Proxy},label={lst:vite-config}]
import { defineConfig } from 'vite';
import react from '@vitejs/plugin-react';

export default defineConfig({
  plugins: [react()],
  server: {
    proxy: {
      '/ws': {
        target: 'ws://localhost:8081',
        ws: true
      }
    }
  },
  build: {
    outDir: 'dist'  // Standard-Output
  }
});
\end{lstlisting}

\begin{infobox}[Konsequenz]
Änderungen im Code erscheinen in unter \SI{100}{\milli\second} im Browser.
Nach \code{npm run build} liegen statische Dateien in \code{dist/} bereit,
die dein Python-Backend ausliefert.
\end{infobox}

% -----------------------------------------------------------------------------
\subsection{ESLint + Prettier -- Code-Qualität automatisieren}
\label{subsec:eslint-prettier}

Inkonsistenter Code (Tabs vs.\ Spaces, fehlende Semikolons) erschwert Reviews.
Potenzielle Bugs bleiben unentdeckt.

\paragraph{ESLint (Analyse)}

ESLint findet problematische Muster:

\begin{lstlisting}[style=shell]
npm install -D eslint @typescript-eslint/parser @typescript-eslint/eslint-plugin
\end{lstlisting}

\begin{lstlisting}[style=base,numbers=left,caption={ESLint-Konfiguration},label={lst:eslint-config}]
// eslint.config.js
export default [
  {
    files: ['**/*.{ts,tsx}'],
    rules: {
      'no-unused-vars': 'error',
      '@typescript-eslint/no-explicit-any': 'warn',
      'react-hooks/exhaustive-deps': 'warn'
    }
  }
];
\end{lstlisting}

\paragraph{Prettier (Formatierung)}

Prettier erzwingt einheitlichen Stil:

\begin{lstlisting}[style=shell]
npm install -D prettier eslint-config-prettier
\end{lstlisting}

\begin{lstlisting}[style=json,caption={Prettier-Konfiguration},label={lst:prettier-config}]
{
  "semi": true,
  "singleQuote": true,
  "tabWidth": 2,
  "trailingComma": "es5"
}
\end{lstlisting}

\paragraph{Integration in VS Code}

\begin{lstlisting}[style=json,caption={VS Code Settings},label={lst:vscode-settings}]
{
  "editor.formatOnSave": true,
  "editor.defaultFormatter": "esbenp.prettier-vscode",
  "editor.codeActionsOnSave": {
    "source.fixAll.eslint": "explicit"
  }
}
\end{lstlisting}

\begin{tipbox}[Konsequenz]
Beim Speichern formatiert Prettier automatisch, ESLint markiert Probleme.
Das eliminiert Style-Diskussionen und fängt häufige Fehler früh ab.
\end{tipbox}

% -----------------------------------------------------------------------------
\subsection{React Router -- Client-seitige Navigation}
\label{subsec:react-router}

Bei klassischen Webseiten lädt jeder Seitenwechsel die komplette Seite neu.
Der WebSocket würde dabei getrennt.

React Router wechselt Komponenten, ohne die Seite neu zu laden:

\begin{lstlisting}[style=shell]
npm install react-router-dom
\end{lstlisting}

\begin{lstlisting}[style=base,numbers=left,caption={React Router Setup},label={lst:router-setup}]
// App.tsx
import { BrowserRouter, Routes, Route, Link } from 'react-router-dom';

function App() {
  return (
    <BrowserRouter>
      <nav>
        <Link to="/">Dashboard</Link>
        <Link to="/config">Konfiguration</Link>
        <Link to="/debug">Debug</Link>
      </nav>

      <Routes>
        <Route path="/" element={<Dashboard />} />
        <Route path="/config" element={<ConfigPanel />} />
        <Route path="/debug" element={<DebugView />} />
      </Routes>
    </BrowserRouter>
  );
}
\end{lstlisting}

\begin{warnbox}[Wichtig für Selection Panel]
Der WebSocket-Hook lebt oberhalb des Routers, sodass die Verbindung bei
Navigation bestehen bleibt:
\end{warnbox}

\begin{lstlisting}[style=base,numbers=left,caption={WebSocket-Context oberhalb des Routers},label={lst:ws-context}]
function App() {
  const ws = useWebSocket('ws://localhost:8081/ws');

  return (
    <WebSocketContext.Provider value={ws}>
      <BrowserRouter>
        {/* Routes hier */}
      </BrowserRouter>
    </WebSocketContext.Provider>
  );
}
\end{lstlisting}

Du kannst zwischen Dashboard, Konfiguration und Debug-Ansicht wechseln,
ohne die ESP32-Verbindung zu verlieren.

% -----------------------------------------------------------------------------
\subsection{Tailwind CSS -- Utility-First Styling}
\label{subsec:tailwind}

Traditionelles CSS erfordert separate Dateien und erfundene Klassennamen:

\begin{lstlisting}[style=base,numbers=none]
/* styles.css - waechst unkontrolliert */
.led-button { ... }
.led-button-active { ... }
.led-button-disabled { ... }
\end{lstlisting}

Tailwind bietet vordefinierte Utility-Klassen direkt im Markup:

\begin{lstlisting}[style=shell]
npm install -D tailwindcss postcss autoprefixer
npx tailwindcss init -p
\end{lstlisting}

\begin{lstlisting}[style=base,numbers=left,caption={Tailwind-Konfiguration},label={lst:tailwind-config}]
// tailwind.config.js
export default {
  content: ['./index.html', './src/**/*.{js,ts,jsx,tsx}'],
  theme: {
    extend: {
      colors: {
        'led-active': '#22c55e',
        'led-inactive': '#374151'
      }
    }
  }
};
\end{lstlisting}

\paragraph{Anwendung}

\begin{lstlisting}[style=base,numbers=left,caption={Tailwind-Klassen in React},label={lst:tailwind-react}]
// Statt CSS-Klassen zu erfinden: direkte Utilities
function LedButton({ isActive }: { isActive: boolean }) {
  return (
    <button className={`
      w-12 h-12 rounded-full
      transition-colors duration-150
      ${isActive
        ? 'bg-led-active shadow-lg shadow-green-500/50'
        : 'bg-led-inactive hover:bg-gray-600'}
    `}>
    </button>
  );
}
\end{lstlisting}

\paragraph{Grid für 100 LEDs}

\begin{lstlisting}[style=base,numbers=none]
<div className="grid grid-cols-10 gap-2 p-4">
  {leds.map((led, i) => <LedButton key={i} isActive={led} />)}
</div>
\end{lstlisting}

\begin{infobox}[Konsequenz]
Kein Wechsel zwischen Dateien, keine CSS-Spezifitätskonflikte. Die Bundle-Größe
bleibt klein, weil Tailwind nur genutzte Klassen einbindet.
\end{infobox}

% -----------------------------------------------------------------------------
\subsection{Zusammenspiel im Selection-Panel-Projekt}
\label{subsec:react-zusammenspiel}

\begin{figure}[H]
\centering
\begin{tikzpicture}[
  box/.style={draw, rectangle, rounded corners=2pt, minimum height=1.2cm,
              font=\ttfamily\footnotesize, fill=white, align=center},
  container/.style={draw, rectangle, rounded corners=3pt, minimum width=13cm,
                    minimum height=4cm, font=\sffamily\small},
  arrow/.style={-{Stealth[length=2mm]}, thick},
  label/.style={font=\sffamily\scriptsize, text=MutedText}
]
  % DEV Container
  \node[container, fill=ArduinoTeal!8] (dev) at (0,4) {};
  \node[font=\sffamily\bfseries, anchor=north west] at (-6.3,5.8) {Entwicklung (Mac)};

  \node[box] (vscode) at (-4,4.2) {VS Code\\+ Extensions};
  \node[box] (vite) at (0,4.2) {Vite Dev\\Server :5173};
  \node[box] (lint) at (4,4.2) {ESLint/\\Prettier};

  \draw[arrow] (vscode) -- (vite);
  \draw[arrow] (lint) -- (vite);

  \node[box] (bridge-dev) at (-4,2.5) {bridge.py\\:8081};
  \node[box, minimum width=2cm] (esp-dev) at (2,2.5) {ESP32-S3};

  \draw[arrow, dashed] (vite.south) -- node[right, label] {proxy /ws} (bridge-dev.north);
  \draw[arrow] (bridge-dev) -- node[above, label] {USB-CDC} (esp-dev);

  % PROD Container
  \node[container, fill=RaspberryRed!8, minimum height=2.5cm] (prod) at (0,-0.5) {};
  \node[font=\sffamily\bfseries, anchor=north west] at (-6.3,0.8) {Produktion (Raspberry Pi)};

  \node[box] (bridge-prod) at (-3,-0.5) {bridge.py\\:8081};
  \node[box, minimum width=2cm] (esp-prod) at (3,-0.5) {ESP32-S3};

  \node[label, anchor=east] at (0,-0.2) {liefert dist/};
  \node[label, anchor=east] at (0,-0.8) {WebSocket /ws};

  \draw[arrow] (bridge-prod) -- node[above, label] {USB-CDC} (esp-prod);
\end{tikzpicture}
\caption{Entwicklungs- und Produktionsumgebung}
\label{fig:dev-prod-setup}
\end{figure}

\paragraph{Deployment-Workflow}

\begin{lstlisting}[style=shell,caption={Frontend-Deployment},label={lst:frontend-deploy}]
# Auf dem Mac
cd apps/dashboard
npm run build          # Erzeugt dist/

# Auf den Pi kopieren
rsync -avh --delete \
  ./apps/dashboard/dist/ \
  pi@rover:/home/pi/embedded_projekt/apps/dashboard/dist/
\end{lstlisting}

% -----------------------------------------------------------------------------
\subsection{Siehe auch}
\label{subsec:react-siehe-auch}

Weiterführende Dokumentation:

\begin{itemize}
  \item \code{apps/dashboard/docs/DEVELOPER.md} -- Frontend-Entwicklung
  \item \code{apps/dashboard/docs/USER\_GUIDE.md} -- Benutzerhandbuch
  \item \code{apps/bridge/docs/PROTOCOL.md} -- WebSocket-Events
  \item \code{apps/bridge/docs/WORKFLOW.md} -- Pi-Deployment
\end{itemize}

% =============================================================================

\section{React-Stack Grundlagen}
\label{sec:react-stack}

% -----------------------------------------------------------------------------
\subsection{Begriffsrahmen und Zielsetzung}
\label{subsec:react-begriffsrahmen}

Du entwickelst das Selection-Panel-Frontend (Browser-UI), das über WebSocket mit
dem Python-Backend kommuniziert. Der Stack ersetzt manuelles HTML/CSS/JS durch
ein wartbares, typsicheres System.

\begin{table}[H]
\centering
\caption{Zentrale Begriffe des React-Stacks}
\label{tab:react-begriffe}
\begin{tabularx}{\textwidth}{@{}l X@{}}
\toprule
\textbf{Begriff} & \textbf{Definition} \\
\midrule
React        & JavaScript-Bibliothek für komponentenbasierte UIs \\
TypeScript   & JavaScript-Erweiterung mit statischer Typisierung \\
Vite         & Build-Tool und Entwicklungsserver (sprich: \glqq viit\grqq, französisch für \glqq schnell\grqq) \\
ESLint       & Statische Code-Analyse (findet Fehler vor Laufzeit) \\
Prettier     & Code-Formatter (einheitlicher Stil) \\
React Router & Client-seitige Navigation ohne Seitenneuladen \\
Tailwind CSS & Utility-First CSS-Framework \\
\bottomrule
\end{tabularx}
\end{table}

% -----------------------------------------------------------------------------
\subsection{React -- Komponenten statt DOM-Manipulation}
\label{subsec:react-komponenten}

Bei klassischem JavaScript manipulierst du das DOM direkt:

\begin{lstlisting}[style=base,numbers=none,caption={Klassische DOM-Manipulation},label={lst:dom-klassisch}]
// Klassisch: imperativ, fehleranfaellig bei Zustandsaenderungen
document.getElementById('led-5').classList.add('active');
document.getElementById('status').textContent = 'LED 5 an';
\end{lstlisting}

Bei 100~LEDs und komplexen Zuständen wird dieser Ansatz unwartbar. Wenn sich ein
Wert ändert, musst du manuell alle betroffenen DOM-Elemente aktualisieren.

React kehrt das Prinzip um -- du beschreibst den \emph{gewünschten Zustand},
React berechnet die nötigen DOM-Änderungen:

\begin{lstlisting}[style=base,numbers=left,caption={React-Komponente (deklarativ)},label={lst:react-komponente}]
// React: deklarativ, Zustand -> UI-Ableitung
function LedButton({ id, isActive, onPress }: LedProps) {
  return (
    <button
      className={isActive ? 'led active' : 'led'}
      onClick={() => onPress(id)}
    >
      LED {id}
    </button>
  );
}
\end{lstlisting}

\paragraph{Kernkonzept -- Komponenten}

Eine Komponente ist eine Funktion, die Props (Eingabeparameter) erhält und JSX
(HTML-ähnliche Syntax) zurückgibt. Komponenten sind wiederverwendbar und isoliert.

\paragraph{Kernkonzept -- State und Hooks}

Lokaler Zustand wird mit \code{useState} verwaltet:

\begin{lstlisting}[style=base,numbers=left,caption={State-Management mit useState},label={lst:usestate}]
function LedPanel() {
  const [activeIds, setActiveIds] = useState<Set<number>>(new Set());

  const handlePress = (id: number) => {
    setActiveIds(prev => {
      const next = new Set(prev);
      next.has(id) ? next.delete(id) : next.add(id);
      return next;
    });
  };

  return (
    <div className="grid">
      {Array.from({ length: 100 }, (_, i) => (
        <LedButton
          key={i}
          id={i}
          isActive={activeIds.has(i)}
          onPress={handlePress}
        />
      ))}
    </div>
  );
}
\end{lstlisting}

\begin{infobox}[Konsequenz für Selection Panel]
Wenn der ESP32 über WebSocket meldet \glqq LED 42 aktiv\grqq, aktualisierst du nur
den State. React rendert automatisch nur die betroffene Komponente neu -- nicht
die anderen 99~LEDs.
\end{infobox}

% -----------------------------------------------------------------------------
\subsection{TypeScript -- Fehler vor Laufzeit finden}
\label{subsec:typescript}

JavaScript-Fehler treten erst zur Laufzeit auf:

\begin{lstlisting}[style=base,numbers=none,caption={JavaScript ohne Typisierung},label={lst:js-untyped}]
// JavaScript: Fehler erst im Browser sichtbar
function setLedState(id, state) {
  websocket.send({ led: id, state: state }); // Typo? Falscher Typ?
}
setLedState("5", "an"); // Kompiliert, crasht zur Laufzeit
\end{lstlisting}

TypeScript prüft Typen beim Kompilieren:

\begin{lstlisting}[style=base,numbers=left,caption={TypeScript mit Typisierung},label={lst:ts-typed}]
// TypeScript: Fehler im Editor, bevor du speicherst
interface LedCommand {
  led: number;
  state: boolean;
}

function setLedState(id: number, state: boolean): void {
  const cmd: LedCommand = { led: id, state };
  websocket.send(JSON.stringify(cmd));
}

setLedState("5", "an"); // Fehler: string ist nicht number
setLedState(5, true);   // korrekt
\end{lstlisting}

\paragraph{Wichtige Typen für dein Projekt}

\begin{lstlisting}[style=base,numbers=left,caption={WebSocket-Nachrichten typisieren},label={lst:ws-types}]
// WebSocket-Nachrichten typisieren
type WsMessage =
  | { type: 'button_press'; id: number }
  | { type: 'led_update'; id: number; state: boolean }
  | { type: 'audio_start'; file: string };

// Props fuer Komponenten
interface LedGridProps {
  leds: Map<number, boolean>;
  onLedClick: (id: number) => void;
}
\end{lstlisting}

\begin{tipbox}[Konsequenz]
Der Compiler fängt Tippfehler, falsche Argumentreihenfolgen und fehlende
Eigenschaften ab. Bei 100~Buttons mit verschiedenen Zuständen reduziert das
die Debugging-Zeit erheblich.
\end{tipbox}

% -----------------------------------------------------------------------------
\subsection{Vite -- Build-Tool und Entwicklungsserver}
\label{subsec:vite}

Browser verstehen weder TypeScript noch JSX direkt. Der Code muss transformiert werden:

\begin{lstlisting}[style=base,numbers=none]
Quellcode (.tsx)  ->  [Transformation]  ->  Browser-Code (.js)
\end{lstlisting}

Ältere Tools wie Webpack benötigen dafür mehrere Sekunden bis Minuten.
Vite nutzt ES-Module (native Browser-Imports) und transformiert nur geänderte Dateien:

\begin{lstlisting}[style=shell,caption={Vite-Projekt erstellen},label={lst:vite-init}]
npm create vite@latest pi-react -- --template react-ts
cd pi-react
npm install
npm run dev   # Startet localhost:5173, Hot Module Replacement
\end{lstlisting}

\paragraph{Zwei Modi}

\begin{table}[H]
\centering
\caption{Vite Build-Modi}
\label{tab:vite-modi}
\begin{tabularx}{\textwidth}{@{}l l X@{}}
\toprule
\textbf{Modus} & \textbf{Befehl} & \textbf{Ergebnis} \\
\midrule
Entwicklung & \code{npm run dev}   & Lokaler Server mit Live-Reload, Source Maps \\
Produktion  & \code{npm run build} & Optimierte statische Dateien in \code{dist/} \\
\bottomrule
\end{tabularx}
\end{table}

\paragraph{Konfiguration (vite.config.ts)}

\begin{lstlisting}[style=base,numbers=left,caption={Vite-Konfiguration mit WebSocket-Proxy},label={lst:vite-config}]
import { defineConfig } from 'vite';
import react from '@vitejs/plugin-react';

export default defineConfig({
  plugins: [react()],
  server: {
    proxy: {
      '/ws': {
        target: 'ws://localhost:8081',
        ws: true
      }
    }
  },
  build: {
    outDir: 'dist'  // Standard-Output
  }
});
\end{lstlisting}

\begin{infobox}[Konsequenz]
Änderungen im Code erscheinen in unter \SI{100}{\milli\second} im Browser.
Nach \code{npm run build} liegen statische Dateien in \code{dist/} bereit,
die dein Python-Backend ausliefert.
\end{infobox}

% -----------------------------------------------------------------------------
\subsection{ESLint + Prettier -- Code-Qualität automatisieren}
\label{subsec:eslint-prettier}

Inkonsistenter Code (Tabs vs.\ Spaces, fehlende Semikolons) erschwert Reviews.
Potenzielle Bugs bleiben unentdeckt.

\paragraph{ESLint (Analyse)}

ESLint findet problematische Muster:

\begin{lstlisting}[style=shell]
npm install -D eslint @typescript-eslint/parser @typescript-eslint/eslint-plugin
\end{lstlisting}

\begin{lstlisting}[style=base,numbers=left,caption={ESLint-Konfiguration},label={lst:eslint-config}]
// eslint.config.js
export default [
  {
    files: ['**/*.{ts,tsx}'],
    rules: {
      'no-unused-vars': 'error',
      '@typescript-eslint/no-explicit-any': 'warn',
      'react-hooks/exhaustive-deps': 'warn'
    }
  }
];
\end{lstlisting}

\paragraph{Prettier (Formatierung)}

Prettier erzwingt einheitlichen Stil:

\begin{lstlisting}[style=shell]
npm install -D prettier eslint-config-prettier
\end{lstlisting}

\begin{lstlisting}[style=json,caption={Prettier-Konfiguration},label={lst:prettier-config}]
{
  "semi": true,
  "singleQuote": true,
  "tabWidth": 2,
  "trailingComma": "es5"
}
\end{lstlisting}

\paragraph{Integration in VS Code}

\begin{lstlisting}[style=json,caption={VS Code Settings},label={lst:vscode-settings}]
{
  "editor.formatOnSave": true,
  "editor.defaultFormatter": "esbenp.prettier-vscode",
  "editor.codeActionsOnSave": {
    "source.fixAll.eslint": "explicit"
  }
}
\end{lstlisting}

\begin{tipbox}[Konsequenz]
Beim Speichern formatiert Prettier automatisch, ESLint markiert Probleme.
Das eliminiert Style-Diskussionen und fängt häufige Fehler früh ab.
\end{tipbox}

% -----------------------------------------------------------------------------
\subsection{React Router -- Client-seitige Navigation}
\label{subsec:react-router}

Bei klassischen Webseiten lädt jeder Seitenwechsel die komplette Seite neu.
Der WebSocket würde dabei getrennt.

React Router wechselt Komponenten, ohne die Seite neu zu laden:

\begin{lstlisting}[style=shell]
npm install react-router-dom
\end{lstlisting}

\begin{lstlisting}[style=base,numbers=left,caption={React Router Setup},label={lst:router-setup}]
// App.tsx
import { BrowserRouter, Routes, Route, Link } from 'react-router-dom';

function App() {
  return (
    <BrowserRouter>
      <nav>
        <Link to="/">Dashboard</Link>
        <Link to="/config">Konfiguration</Link>
        <Link to="/debug">Debug</Link>
      </nav>

      <Routes>
        <Route path="/" element={<Dashboard />} />
        <Route path="/config" element={<ConfigPanel />} />
        <Route path="/debug" element={<DebugView />} />
      </Routes>
    </BrowserRouter>
  );
}
\end{lstlisting}

\begin{warnbox}[Wichtig für Selection Panel]
Der WebSocket-Hook lebt oberhalb des Routers, sodass die Verbindung bei
Navigation bestehen bleibt:
\end{warnbox}

\begin{lstlisting}[style=base,numbers=left,caption={WebSocket-Context oberhalb des Routers},label={lst:ws-context}]
function App() {
  const ws = useWebSocket('ws://localhost:8081/ws');

  return (
    <WebSocketContext.Provider value={ws}>
      <BrowserRouter>
        {/* Routes hier */}
      </BrowserRouter>
    </WebSocketContext.Provider>
  );
}
\end{lstlisting}

Du kannst zwischen Dashboard, Konfiguration und Debug-Ansicht wechseln,
ohne die ESP32-Verbindung zu verlieren.

% -----------------------------------------------------------------------------
\subsection{Tailwind CSS -- Utility-First Styling}
\label{subsec:tailwind}

Traditionelles CSS erfordert separate Dateien und erfundene Klassennamen:

\begin{lstlisting}[style=base,numbers=none]
/* styles.css - waechst unkontrolliert */
.led-button { ... }
.led-button-active { ... }
.led-button-disabled { ... }
\end{lstlisting}

Tailwind bietet vordefinierte Utility-Klassen direkt im Markup:

\begin{lstlisting}[style=shell]
npm install -D tailwindcss postcss autoprefixer
npx tailwindcss init -p
\end{lstlisting}

\begin{lstlisting}[style=base,numbers=left,caption={Tailwind-Konfiguration},label={lst:tailwind-config}]
// tailwind.config.js
export default {
  content: ['./index.html', './src/**/*.{js,ts,jsx,tsx}'],
  theme: {
    extend: {
      colors: {
        'led-active': '#22c55e',
        'led-inactive': '#374151'
      }
    }
  }
};
\end{lstlisting}

\paragraph{Anwendung}

\begin{lstlisting}[style=base,numbers=left,caption={Tailwind-Klassen in React},label={lst:tailwind-react}]
// Statt CSS-Klassen zu erfinden: direkte Utilities
function LedButton({ isActive }: { isActive: boolean }) {
  return (
    <button className={`
      w-12 h-12 rounded-full
      transition-colors duration-150
      ${isActive
        ? 'bg-led-active shadow-lg shadow-green-500/50'
        : 'bg-led-inactive hover:bg-gray-600'}
    `}>
    </button>
  );
}
\end{lstlisting}

\paragraph{Grid für 100 LEDs}

\begin{lstlisting}[style=base,numbers=none]
<div className="grid grid-cols-10 gap-2 p-4">
  {leds.map((led, i) => <LedButton key={i} isActive={led} />)}
</div>
\end{lstlisting}

\begin{infobox}[Konsequenz]
Kein Wechsel zwischen Dateien, keine CSS-Spezifitätskonflikte. Die Bundle-Größe
bleibt klein, weil Tailwind nur genutzte Klassen einbindet.
\end{infobox}

% -----------------------------------------------------------------------------
\subsection{Zusammenspiel im Selection-Panel-Projekt}
\label{subsec:react-zusammenspiel}

\begin{figure}[H]
\centering
\begin{tikzpicture}[
  box/.style={draw, rectangle, rounded corners=2pt, minimum height=1.2cm,
              font=\ttfamily\footnotesize, fill=white, align=center},
  container/.style={draw, rectangle, rounded corners=3pt, minimum width=13cm,
                    minimum height=4cm, font=\sffamily\small},
  arrow/.style={-{Stealth[length=2mm]}, thick},
  label/.style={font=\sffamily\scriptsize, text=MutedText}
]
  % DEV Container
  \node[container, fill=ArduinoTeal!8] (dev) at (0,4) {};
  \node[font=\sffamily\bfseries, anchor=north west] at (-6.3,5.8) {Entwicklung (Mac)};

  \node[box] (vscode) at (-4,4.2) {VS Code\\+ Extensions};
  \node[box] (vite) at (0,4.2) {Vite Dev\\Server :5173};
  \node[box] (lint) at (4,4.2) {ESLint/\\Prettier};

  \draw[arrow] (vscode) -- (vite);
  \draw[arrow] (lint) -- (vite);

  \node[box] (bridge-dev) at (-4,2.5) {bridge.py\\:8081};
  \node[box, minimum width=2cm] (esp-dev) at (2,2.5) {ESP32-S3};

  \draw[arrow, dashed] (vite.south) -- node[right, label] {proxy /ws} (bridge-dev.north);
  \draw[arrow] (bridge-dev) -- node[above, label] {USB-CDC} (esp-dev);

  % PROD Container
  \node[container, fill=RaspberryRed!8, minimum height=2.5cm] (prod) at (0,-0.5) {};
  \node[font=\sffamily\bfseries, anchor=north west] at (-6.3,0.8) {Produktion (Raspberry Pi)};

  \node[box] (bridge-prod) at (-3,-0.5) {bridge.py\\:8081};
  \node[box, minimum width=2cm] (esp-prod) at (3,-0.5) {ESP32-S3};

  \node[label, anchor=east] at (0,-0.2) {liefert dist/};
  \node[label, anchor=east] at (0,-0.8) {WebSocket /ws};

  \draw[arrow] (bridge-prod) -- node[above, label] {USB-CDC} (esp-prod);
\end{tikzpicture}
\caption{Entwicklungs- und Produktionsumgebung}
\label{fig:dev-prod-setup}
\end{figure}

\paragraph{Deployment-Workflow}

\begin{lstlisting}[style=shell,caption={Frontend-Deployment},label={lst:frontend-deploy}]
# Auf dem Mac
cd apps/dashboard
npm run build          # Erzeugt dist/

# Auf den Pi kopieren
rsync -avh --delete \
  ./apps/dashboard/dist/ \
  pi@rover:/home/pi/embedded_projekt/apps/dashboard/dist/
\end{lstlisting}

% -----------------------------------------------------------------------------
\subsection{Siehe auch}
\label{subsec:react-siehe-auch}

Weiterführende Dokumentation:

\begin{itemize}
  \item \code{apps/dashboard/docs/DEVELOPER.md} -- Frontend-Entwicklung
  \item \code{apps/dashboard/docs/USER\_GUIDE.md} -- Benutzerhandbuch
  \item \code{apps/bridge/docs/PROTOCOL.md} -- WebSocket-Events
  \item \code{apps/bridge/docs/WORKFLOW.md} -- Pi-Deployment
\end{itemize}

% =============================================================================

\section{React-Stack Grundlagen}
\label{sec:react-stack}

% -----------------------------------------------------------------------------
\subsection{Begriffsrahmen und Zielsetzung}
\label{subsec:react-begriffsrahmen}

Du entwickelst das Selection-Panel-Frontend (Browser-UI), das über WebSocket mit
dem Python-Backend kommuniziert. Der Stack ersetzt manuelles HTML/CSS/JS durch
ein wartbares, typsicheres System.

\begin{table}[H]
\centering
\caption{Zentrale Begriffe des React-Stacks}
\label{tab:react-begriffe}
\begin{tabularx}{\textwidth}{@{}l X@{}}
\toprule
\textbf{Begriff} & \textbf{Definition} \\
\midrule
React        & JavaScript-Bibliothek für komponentenbasierte UIs \\
TypeScript   & JavaScript-Erweiterung mit statischer Typisierung \\
Vite         & Build-Tool und Entwicklungsserver (sprich: \glqq viit\grqq, französisch für \glqq schnell\grqq) \\
ESLint       & Statische Code-Analyse (findet Fehler vor Laufzeit) \\
Prettier     & Code-Formatter (einheitlicher Stil) \\
React Router & Client-seitige Navigation ohne Seitenneuladen \\
Tailwind CSS & Utility-First CSS-Framework \\
\bottomrule
\end{tabularx}
\end{table}

% -----------------------------------------------------------------------------
\subsection{React -- Komponenten statt DOM-Manipulation}
\label{subsec:react-komponenten}

Bei klassischem JavaScript manipulierst du das DOM direkt:

\begin{lstlisting}[style=base,numbers=none,caption={Klassische DOM-Manipulation},label={lst:dom-klassisch}]
// Klassisch: imperativ, fehleranfaellig bei Zustandsaenderungen
document.getElementById('led-5').classList.add('active');
document.getElementById('status').textContent = 'LED 5 an';
\end{lstlisting}

Bei 100~LEDs und komplexen Zuständen wird dieser Ansatz unwartbar. Wenn sich ein
Wert ändert, musst du manuell alle betroffenen DOM-Elemente aktualisieren.

React kehrt das Prinzip um -- du beschreibst den \emph{gewünschten Zustand},
React berechnet die nötigen DOM-Änderungen:

\begin{lstlisting}[style=base,numbers=left,caption={React-Komponente (deklarativ)},label={lst:react-komponente}]
// React: deklarativ, Zustand -> UI-Ableitung
function LedButton({ id, isActive, onPress }: LedProps) {
  return (
    <button
      className={isActive ? 'led active' : 'led'}
      onClick={() => onPress(id)}
    >
      LED {id}
    </button>
  );
}
\end{lstlisting}

\paragraph{Kernkonzept -- Komponenten}

Eine Komponente ist eine Funktion, die Props (Eingabeparameter) erhält und JSX
(HTML-ähnliche Syntax) zurückgibt. Komponenten sind wiederverwendbar und isoliert.

\paragraph{Kernkonzept -- State und Hooks}

Lokaler Zustand wird mit \code{useState} verwaltet:

\begin{lstlisting}[style=base,numbers=left,caption={State-Management mit useState},label={lst:usestate}]
function LedPanel() {
  const [activeIds, setActiveIds] = useState<Set<number>>(new Set());

  const handlePress = (id: number) => {
    setActiveIds(prev => {
      const next = new Set(prev);
      next.has(id) ? next.delete(id) : next.add(id);
      return next;
    });
  };

  return (
    <div className="grid">
      {Array.from({ length: 100 }, (_, i) => (
        <LedButton
          key={i}
          id={i}
          isActive={activeIds.has(i)}
          onPress={handlePress}
        />
      ))}
    </div>
  );
}
\end{lstlisting}

\begin{infobox}[Konsequenz für Selection Panel]
Wenn der ESP32 über WebSocket meldet \glqq LED 42 aktiv\grqq, aktualisierst du nur
den State. React rendert automatisch nur die betroffene Komponente neu -- nicht
die anderen 99~LEDs.
\end{infobox}

% -----------------------------------------------------------------------------
\subsection{TypeScript -- Fehler vor Laufzeit finden}
\label{subsec:typescript}

JavaScript-Fehler treten erst zur Laufzeit auf:

\begin{lstlisting}[style=base,numbers=none,caption={JavaScript ohne Typisierung},label={lst:js-untyped}]
// JavaScript: Fehler erst im Browser sichtbar
function setLedState(id, state) {
  websocket.send({ led: id, state: state }); // Typo? Falscher Typ?
}
setLedState("5", "an"); // Kompiliert, crasht zur Laufzeit
\end{lstlisting}

TypeScript prüft Typen beim Kompilieren:

\begin{lstlisting}[style=base,numbers=left,caption={TypeScript mit Typisierung},label={lst:ts-typed}]
// TypeScript: Fehler im Editor, bevor du speicherst
interface LedCommand {
  led: number;
  state: boolean;
}

function setLedState(id: number, state: boolean): void {
  const cmd: LedCommand = { led: id, state };
  websocket.send(JSON.stringify(cmd));
}

setLedState("5", "an"); // Fehler: string ist nicht number
setLedState(5, true);   // korrekt
\end{lstlisting}

\paragraph{Wichtige Typen für dein Projekt}

\begin{lstlisting}[style=base,numbers=left,caption={WebSocket-Nachrichten typisieren},label={lst:ws-types}]
// WebSocket-Nachrichten typisieren
type WsMessage =
  | { type: 'button_press'; id: number }
  | { type: 'led_update'; id: number; state: boolean }
  | { type: 'audio_start'; file: string };

// Props fuer Komponenten
interface LedGridProps {
  leds: Map<number, boolean>;
  onLedClick: (id: number) => void;
}
\end{lstlisting}

\begin{tipbox}[Konsequenz]
Der Compiler fängt Tippfehler, falsche Argumentreihenfolgen und fehlende
Eigenschaften ab. Bei 100~Buttons mit verschiedenen Zuständen reduziert das
die Debugging-Zeit erheblich.
\end{tipbox}

% -----------------------------------------------------------------------------
\subsection{Vite -- Build-Tool und Entwicklungsserver}
\label{subsec:vite}

Browser verstehen weder TypeScript noch JSX direkt. Der Code muss transformiert werden:

\begin{lstlisting}[style=base,numbers=none]
Quellcode (.tsx)  ->  [Transformation]  ->  Browser-Code (.js)
\end{lstlisting}

Ältere Tools wie Webpack benötigen dafür mehrere Sekunden bis Minuten.
Vite nutzt ES-Module (native Browser-Imports) und transformiert nur geänderte Dateien:

\begin{lstlisting}[style=shell,caption={Vite-Projekt erstellen},label={lst:vite-init}]
npm create vite@latest pi-react -- --template react-ts
cd pi-react
npm install
npm run dev   # Startet localhost:5173, Hot Module Replacement
\end{lstlisting}

\paragraph{Zwei Modi}

\begin{table}[H]
\centering
\caption{Vite Build-Modi}
\label{tab:vite-modi}
\begin{tabularx}{\textwidth}{@{}l l X@{}}
\toprule
\textbf{Modus} & \textbf{Befehl} & \textbf{Ergebnis} \\
\midrule
Entwicklung & \code{npm run dev}   & Lokaler Server mit Live-Reload, Source Maps \\
Produktion  & \code{npm run build} & Optimierte statische Dateien in \code{dist/} \\
\bottomrule
\end{tabularx}
\end{table}

\paragraph{Konfiguration (vite.config.ts)}

\begin{lstlisting}[style=base,numbers=left,caption={Vite-Konfiguration mit WebSocket-Proxy},label={lst:vite-config}]
import { defineConfig } from 'vite';
import react from '@vitejs/plugin-react';

export default defineConfig({
  plugins: [react()],
  server: {
    proxy: {
      '/ws': {
        target: 'ws://localhost:8081',
        ws: true
      }
    }
  },
  build: {
    outDir: 'dist'  // Standard-Output
  }
});
\end{lstlisting}

\begin{infobox}[Konsequenz]
Änderungen im Code erscheinen in unter \SI{100}{\milli\second} im Browser.
Nach \code{npm run build} liegen statische Dateien in \code{dist/} bereit,
die dein Python-Backend ausliefert.
\end{infobox}

% -----------------------------------------------------------------------------
\subsection{ESLint + Prettier -- Code-Qualität automatisieren}
\label{subsec:eslint-prettier}

Inkonsistenter Code (Tabs vs.\ Spaces, fehlende Semikolons) erschwert Reviews.
Potenzielle Bugs bleiben unentdeckt.

\paragraph{ESLint (Analyse)}

ESLint findet problematische Muster:

\begin{lstlisting}[style=shell]
npm install -D eslint @typescript-eslint/parser @typescript-eslint/eslint-plugin
\end{lstlisting}

\begin{lstlisting}[style=base,numbers=left,caption={ESLint-Konfiguration},label={lst:eslint-config}]
// eslint.config.js
export default [
  {
    files: ['**/*.{ts,tsx}'],
    rules: {
      'no-unused-vars': 'error',
      '@typescript-eslint/no-explicit-any': 'warn',
      'react-hooks/exhaustive-deps': 'warn'
    }
  }
];
\end{lstlisting}

\paragraph{Prettier (Formatierung)}

Prettier erzwingt einheitlichen Stil:

\begin{lstlisting}[style=shell]
npm install -D prettier eslint-config-prettier
\end{lstlisting}

\begin{lstlisting}[style=json,caption={Prettier-Konfiguration},label={lst:prettier-config}]
{
  "semi": true,
  "singleQuote": true,
  "tabWidth": 2,
  "trailingComma": "es5"
}
\end{lstlisting}

\paragraph{Integration in VS Code}

\begin{lstlisting}[style=json,caption={VS Code Settings},label={lst:vscode-settings}]
{
  "editor.formatOnSave": true,
  "editor.defaultFormatter": "esbenp.prettier-vscode",
  "editor.codeActionsOnSave": {
    "source.fixAll.eslint": "explicit"
  }
}
\end{lstlisting}

\begin{tipbox}[Konsequenz]
Beim Speichern formatiert Prettier automatisch, ESLint markiert Probleme.
Das eliminiert Style-Diskussionen und fängt häufige Fehler früh ab.
\end{tipbox}

% -----------------------------------------------------------------------------
\subsection{React Router -- Client-seitige Navigation}
\label{subsec:react-router}

Bei klassischen Webseiten lädt jeder Seitenwechsel die komplette Seite neu.
Der WebSocket würde dabei getrennt.

React Router wechselt Komponenten, ohne die Seite neu zu laden:

\begin{lstlisting}[style=shell]
npm install react-router-dom
\end{lstlisting}

\begin{lstlisting}[style=base,numbers=left,caption={React Router Setup},label={lst:router-setup}]
// App.tsx
import { BrowserRouter, Routes, Route, Link } from 'react-router-dom';

function App() {
  return (
    <BrowserRouter>
      <nav>
        <Link to="/">Dashboard</Link>
        <Link to="/config">Konfiguration</Link>
        <Link to="/debug">Debug</Link>
      </nav>

      <Routes>
        <Route path="/" element={<Dashboard />} />
        <Route path="/config" element={<ConfigPanel />} />
        <Route path="/debug" element={<DebugView />} />
      </Routes>
    </BrowserRouter>
  );
}
\end{lstlisting}

\begin{warnbox}[Wichtig für Selection Panel]
Der WebSocket-Hook lebt oberhalb des Routers, sodass die Verbindung bei
Navigation bestehen bleibt:
\end{warnbox}

\begin{lstlisting}[style=base,numbers=left,caption={WebSocket-Context oberhalb des Routers},label={lst:ws-context}]
function App() {
  const ws = useWebSocket('ws://localhost:8081/ws');

  return (
    <WebSocketContext.Provider value={ws}>
      <BrowserRouter>
        {/* Routes hier */}
      </BrowserRouter>
    </WebSocketContext.Provider>
  );
}
\end{lstlisting}

Du kannst zwischen Dashboard, Konfiguration und Debug-Ansicht wechseln,
ohne die ESP32-Verbindung zu verlieren.

% -----------------------------------------------------------------------------
\subsection{Tailwind CSS -- Utility-First Styling}
\label{subsec:tailwind}

Traditionelles CSS erfordert separate Dateien und erfundene Klassennamen:

\begin{lstlisting}[style=base,numbers=none]
/* styles.css - waechst unkontrolliert */
.led-button { ... }
.led-button-active { ... }
.led-button-disabled { ... }
\end{lstlisting}

Tailwind bietet vordefinierte Utility-Klassen direkt im Markup:

\begin{lstlisting}[style=shell]
npm install -D tailwindcss postcss autoprefixer
npx tailwindcss init -p
\end{lstlisting}

\begin{lstlisting}[style=base,numbers=left,caption={Tailwind-Konfiguration},label={lst:tailwind-config}]
// tailwind.config.js
export default {
  content: ['./index.html', './src/**/*.{js,ts,jsx,tsx}'],
  theme: {
    extend: {
      colors: {
        'led-active': '#22c55e',
        'led-inactive': '#374151'
      }
    }
  }
};
\end{lstlisting}

\paragraph{Anwendung}

\begin{lstlisting}[style=base,numbers=left,caption={Tailwind-Klassen in React},label={lst:tailwind-react}]
// Statt CSS-Klassen zu erfinden: direkte Utilities
function LedButton({ isActive }: { isActive: boolean }) {
  return (
    <button className={`
      w-12 h-12 rounded-full
      transition-colors duration-150
      ${isActive
        ? 'bg-led-active shadow-lg shadow-green-500/50'
        : 'bg-led-inactive hover:bg-gray-600'}
    `}>
    </button>
  );
}
\end{lstlisting}

\paragraph{Grid für 100 LEDs}

\begin{lstlisting}[style=base,numbers=none]
<div className="grid grid-cols-10 gap-2 p-4">
  {leds.map((led, i) => <LedButton key={i} isActive={led} />)}
</div>
\end{lstlisting}

\begin{infobox}[Konsequenz]
Kein Wechsel zwischen Dateien, keine CSS-Spezifitätskonflikte. Die Bundle-Größe
bleibt klein, weil Tailwind nur genutzte Klassen einbindet.
\end{infobox}

% -----------------------------------------------------------------------------
\subsection{Zusammenspiel im Selection-Panel-Projekt}
\label{subsec:react-zusammenspiel}

\begin{figure}[H]
\centering
\begin{tikzpicture}[
  box/.style={draw, rectangle, rounded corners=2pt, minimum height=1.2cm,
              font=\ttfamily\footnotesize, fill=white, align=center},
  container/.style={draw, rectangle, rounded corners=3pt, minimum width=13cm,
                    minimum height=4cm, font=\sffamily\small},
  arrow/.style={-{Stealth[length=2mm]}, thick},
  label/.style={font=\sffamily\scriptsize, text=MutedText}
]
  % DEV Container
  \node[container, fill=ArduinoTeal!8] (dev) at (0,4) {};
  \node[font=\sffamily\bfseries, anchor=north west] at (-6.3,5.8) {Entwicklung (Mac)};

  \node[box] (vscode) at (-4,4.2) {VS Code\\+ Extensions};
  \node[box] (vite) at (0,4.2) {Vite Dev\\Server :5173};
  \node[box] (lint) at (4,4.2) {ESLint/\\Prettier};

  \draw[arrow] (vscode) -- (vite);
  \draw[arrow] (lint) -- (vite);

  \node[box] (bridge-dev) at (-4,2.5) {bridge.py\\:8081};
  \node[box, minimum width=2cm] (esp-dev) at (2,2.5) {ESP32-S3};

  \draw[arrow, dashed] (vite.south) -- node[right, label] {proxy /ws} (bridge-dev.north);
  \draw[arrow] (bridge-dev) -- node[above, label] {USB-CDC} (esp-dev);

  % PROD Container
  \node[container, fill=RaspberryRed!8, minimum height=2.5cm] (prod) at (0,-0.5) {};
  \node[font=\sffamily\bfseries, anchor=north west] at (-6.3,0.8) {Produktion (Raspberry Pi)};

  \node[box] (bridge-prod) at (-3,-0.5) {bridge.py\\:8081};
  \node[box, minimum width=2cm] (esp-prod) at (3,-0.5) {ESP32-S3};

  \node[label, anchor=east] at (0,-0.2) {liefert dist/};
  \node[label, anchor=east] at (0,-0.8) {WebSocket /ws};

  \draw[arrow] (bridge-prod) -- node[above, label] {USB-CDC} (esp-prod);
\end{tikzpicture}
\caption{Entwicklungs- und Produktionsumgebung}
\label{fig:dev-prod-setup}
\end{figure}

\paragraph{Deployment-Workflow}

\begin{lstlisting}[style=shell,caption={Frontend-Deployment},label={lst:frontend-deploy}]
# Auf dem Mac
cd apps/dashboard
npm run build          # Erzeugt dist/

# Auf den Pi kopieren
rsync -avh --delete \
  ./apps/dashboard/dist/ \
  pi@rover:/home/pi/embedded_projekt/apps/dashboard/dist/
\end{lstlisting}

% -----------------------------------------------------------------------------
\subsection{Siehe auch}
\label{subsec:react-siehe-auch}

Weiterführende Dokumentation:

\begin{itemize}
  \item \code{apps/dashboard/docs/DEVELOPER.md} -- Frontend-Entwicklung
  \item \code{apps/dashboard/docs/USER\_GUIDE.md} -- Benutzerhandbuch
  \item \code{apps/bridge/docs/PROTOCOL.md} -- WebSocket-Events
  \item \code{apps/bridge/docs/WORKFLOW.md} -- Pi-Deployment
\end{itemize}

% =============================================================================

\section{React-Stack Grundlagen}
\label{sec:react-stack}

% -----------------------------------------------------------------------------
\subsection{Begriffsrahmen und Zielsetzung}
\label{subsec:react-begriffsrahmen}

Du entwickelst das Selection-Panel-Frontend (Browser-UI), das über WebSocket mit
dem Python-Backend kommuniziert. Der Stack ersetzt manuelles HTML/CSS/JS durch
ein wartbares, typsicheres System.

\begin{table}[H]
\centering
\caption{Zentrale Begriffe des React-Stacks}
\label{tab:react-begriffe}
\begin{tabularx}{\textwidth}{@{}l X@{}}
\toprule
\textbf{Begriff} & \textbf{Definition} \\
\midrule
React        & JavaScript-Bibliothek für komponentenbasierte UIs \\
TypeScript   & JavaScript-Erweiterung mit statischer Typisierung \\
Vite         & Build-Tool und Entwicklungsserver (sprich: \glqq viit\grqq, französisch für \glqq schnell\grqq) \\
ESLint       & Statische Code-Analyse (findet Fehler vor Laufzeit) \\
Prettier     & Code-Formatter (einheitlicher Stil) \\
React Router & Client-seitige Navigation ohne Seitenneuladen \\
Tailwind CSS & Utility-First CSS-Framework \\
\bottomrule
\end{tabularx}
\end{table}

% -----------------------------------------------------------------------------
\subsection{React -- Komponenten statt DOM-Manipulation}
\label{subsec:react-komponenten}

Bei klassischem JavaScript manipulierst du das DOM direkt:

\begin{lstlisting}[style=base,numbers=none,caption={Klassische DOM-Manipulation},label={lst:dom-klassisch}]
// Klassisch: imperativ, fehleranfaellig bei Zustandsaenderungen
document.getElementById('led-5').classList.add('active');
document.getElementById('status').textContent = 'LED 5 an';
\end{lstlisting}

Bei 100~LEDs und komplexen Zuständen wird dieser Ansatz unwartbar. Wenn sich ein
Wert ändert, musst du manuell alle betroffenen DOM-Elemente aktualisieren.

React kehrt das Prinzip um -- du beschreibst den \emph{gewünschten Zustand},
React berechnet die nötigen DOM-Änderungen:

\begin{lstlisting}[style=base,numbers=left,caption={React-Komponente (deklarativ)},label={lst:react-komponente}]
// React: deklarativ, Zustand -> UI-Ableitung
function LedButton({ id, isActive, onPress }: LedProps) {
  return (
    <button
      className={isActive ? 'led active' : 'led'}
      onClick={() => onPress(id)}
    >
      LED {id}
    </button>
  );
}
\end{lstlisting}

\paragraph{Kernkonzept -- Komponenten}

Eine Komponente ist eine Funktion, die Props (Eingabeparameter) erhält und JSX
(HTML-ähnliche Syntax) zurückgibt. Komponenten sind wiederverwendbar und isoliert.

\paragraph{Kernkonzept -- State und Hooks}

Lokaler Zustand wird mit \code{useState} verwaltet:

\begin{lstlisting}[style=base,numbers=left,caption={State-Management mit useState},label={lst:usestate}]
function LedPanel() {
  const [activeIds, setActiveIds] = useState<Set<number>>(new Set());

  const handlePress = (id: number) => {
    setActiveIds(prev => {
      const next = new Set(prev);
      next.has(id) ? next.delete(id) : next.add(id);
      return next;
    });
  };

  return (
    <div className="grid">
      {Array.from({ length: 100 }, (_, i) => (
        <LedButton
          key={i}
          id={i}
          isActive={activeIds.has(i)}
          onPress={handlePress}
        />
      ))}
    </div>
  );
}
\end{lstlisting}

\begin{infobox}[Konsequenz für Selection Panel]
Wenn der ESP32 über WebSocket meldet \glqq LED 42 aktiv\grqq, aktualisierst du nur
den State. React rendert automatisch nur die betroffene Komponente neu -- nicht
die anderen 99~LEDs.
\end{infobox}

% -----------------------------------------------------------------------------
\subsection{TypeScript -- Fehler vor Laufzeit finden}
\label{subsec:typescript}

JavaScript-Fehler treten erst zur Laufzeit auf:

\begin{lstlisting}[style=base,numbers=none,caption={JavaScript ohne Typisierung},label={lst:js-untyped}]
// JavaScript: Fehler erst im Browser sichtbar
function setLedState(id, state) {
  websocket.send({ led: id, state: state }); // Typo? Falscher Typ?
}
setLedState("5", "an"); // Kompiliert, crasht zur Laufzeit
\end{lstlisting}

TypeScript prüft Typen beim Kompilieren:

\begin{lstlisting}[style=base,numbers=left,caption={TypeScript mit Typisierung},label={lst:ts-typed}]
// TypeScript: Fehler im Editor, bevor du speicherst
interface LedCommand {
  led: number;
  state: boolean;
}

function setLedState(id: number, state: boolean): void {
  const cmd: LedCommand = { led: id, state };
  websocket.send(JSON.stringify(cmd));
}

setLedState("5", "an"); // Fehler: string ist nicht number
setLedState(5, true);   // korrekt
\end{lstlisting}

\paragraph{Wichtige Typen für dein Projekt}

\begin{lstlisting}[style=base,numbers=left,caption={WebSocket-Nachrichten typisieren},label={lst:ws-types}]
// WebSocket-Nachrichten typisieren
type WsMessage =
  | { type: 'button_press'; id: number }
  | { type: 'led_update'; id: number; state: boolean }
  | { type: 'audio_start'; file: string };

// Props fuer Komponenten
interface LedGridProps {
  leds: Map<number, boolean>;
  onLedClick: (id: number) => void;
}
\end{lstlisting}

\begin{tipbox}[Konsequenz]
Der Compiler fängt Tippfehler, falsche Argumentreihenfolgen und fehlende
Eigenschaften ab. Bei 100~Buttons mit verschiedenen Zuständen reduziert das
die Debugging-Zeit erheblich.
\end{tipbox}

% -----------------------------------------------------------------------------
\subsection{Vite -- Build-Tool und Entwicklungsserver}
\label{subsec:vite}

Browser verstehen weder TypeScript noch JSX direkt. Der Code muss transformiert werden:

\begin{lstlisting}[style=base,numbers=none]
Quellcode (.tsx)  ->  [Transformation]  ->  Browser-Code (.js)
\end{lstlisting}

Ältere Tools wie Webpack benötigen dafür mehrere Sekunden bis Minuten.
Vite nutzt ES-Module (native Browser-Imports) und transformiert nur geänderte Dateien:

\begin{lstlisting}[style=shell,caption={Vite-Projekt erstellen},label={lst:vite-init}]
npm create vite@latest pi-react -- --template react-ts
cd pi-react
npm install
npm run dev   # Startet localhost:5173, Hot Module Replacement
\end{lstlisting}

\paragraph{Zwei Modi}

\begin{table}[H]
\centering
\caption{Vite Build-Modi}
\label{tab:vite-modi}
\begin{tabularx}{\textwidth}{@{}l l X@{}}
\toprule
\textbf{Modus} & \textbf{Befehl} & \textbf{Ergebnis} \\
\midrule
Entwicklung & \code{npm run dev}   & Lokaler Server mit Live-Reload, Source Maps \\
Produktion  & \code{npm run build} & Optimierte statische Dateien in \code{dist/} \\
\bottomrule
\end{tabularx}
\end{table}

\paragraph{Konfiguration (vite.config.ts)}

\begin{lstlisting}[style=base,numbers=left,caption={Vite-Konfiguration mit WebSocket-Proxy},label={lst:vite-config}]
import { defineConfig } from 'vite';
import react from '@vitejs/plugin-react';

export default defineConfig({
  plugins: [react()],
  server: {
    proxy: {
      '/ws': {
        target: 'ws://localhost:8081',
        ws: true
      }
    }
  },
  build: {
    outDir: 'dist'  // Standard-Output
  }
});
\end{lstlisting}

\begin{infobox}[Konsequenz]
Änderungen im Code erscheinen in unter \SI{100}{\milli\second} im Browser.
Nach \code{npm run build} liegen statische Dateien in \code{dist/} bereit,
die dein Python-Backend ausliefert.
\end{infobox}

% -----------------------------------------------------------------------------
\subsection{ESLint + Prettier -- Code-Qualität automatisieren}
\label{subsec:eslint-prettier}

Inkonsistenter Code (Tabs vs.\ Spaces, fehlende Semikolons) erschwert Reviews.
Potenzielle Bugs bleiben unentdeckt.

\paragraph{ESLint (Analyse)}

ESLint findet problematische Muster:

\begin{lstlisting}[style=shell]
npm install -D eslint @typescript-eslint/parser @typescript-eslint/eslint-plugin
\end{lstlisting}

\begin{lstlisting}[style=base,numbers=left,caption={ESLint-Konfiguration},label={lst:eslint-config}]
// eslint.config.js
export default [
  {
    files: ['**/*.{ts,tsx}'],
    rules: {
      'no-unused-vars': 'error',
      '@typescript-eslint/no-explicit-any': 'warn',
      'react-hooks/exhaustive-deps': 'warn'
    }
  }
];
\end{lstlisting}

\paragraph{Prettier (Formatierung)}

Prettier erzwingt einheitlichen Stil:

\begin{lstlisting}[style=shell]
npm install -D prettier eslint-config-prettier
\end{lstlisting}

\begin{lstlisting}[style=json,caption={Prettier-Konfiguration},label={lst:prettier-config}]
{
  "semi": true,
  "singleQuote": true,
  "tabWidth": 2,
  "trailingComma": "es5"
}
\end{lstlisting}

\paragraph{Integration in VS Code}

\begin{lstlisting}[style=json,caption={VS Code Settings},label={lst:vscode-settings}]
{
  "editor.formatOnSave": true,
  "editor.defaultFormatter": "esbenp.prettier-vscode",
  "editor.codeActionsOnSave": {
    "source.fixAll.eslint": "explicit"
  }
}
\end{lstlisting}

\begin{tipbox}[Konsequenz]
Beim Speichern formatiert Prettier automatisch, ESLint markiert Probleme.
Das eliminiert Style-Diskussionen und fängt häufige Fehler früh ab.
\end{tipbox}

% -----------------------------------------------------------------------------
\subsection{React Router -- Client-seitige Navigation}
\label{subsec:react-router}

Bei klassischen Webseiten lädt jeder Seitenwechsel die komplette Seite neu.
Der WebSocket würde dabei getrennt.

React Router wechselt Komponenten, ohne die Seite neu zu laden:

\begin{lstlisting}[style=shell]
npm install react-router-dom
\end{lstlisting}

\begin{lstlisting}[style=base,numbers=left,caption={React Router Setup},label={lst:router-setup}]
// App.tsx
import { BrowserRouter, Routes, Route, Link } from 'react-router-dom';

function App() {
  return (
    <BrowserRouter>
      <nav>
        <Link to="/">Dashboard</Link>
        <Link to="/config">Konfiguration</Link>
        <Link to="/debug">Debug</Link>
      </nav>

      <Routes>
        <Route path="/" element={<Dashboard />} />
        <Route path="/config" element={<ConfigPanel />} />
        <Route path="/debug" element={<DebugView />} />
      </Routes>
    </BrowserRouter>
  );
}
\end{lstlisting}

\begin{warnbox}[Wichtig für Selection Panel]
Der WebSocket-Hook lebt oberhalb des Routers, sodass die Verbindung bei
Navigation bestehen bleibt:
\end{warnbox}

\begin{lstlisting}[style=base,numbers=left,caption={WebSocket-Context oberhalb des Routers},label={lst:ws-context}]
function App() {
  const ws = useWebSocket('ws://localhost:8081/ws');

  return (
    <WebSocketContext.Provider value={ws}>
      <BrowserRouter>
        {/* Routes hier */}
      </BrowserRouter>
    </WebSocketContext.Provider>
  );
}
\end{lstlisting}

Du kannst zwischen Dashboard, Konfiguration und Debug-Ansicht wechseln,
ohne die ESP32-Verbindung zu verlieren.

% -----------------------------------------------------------------------------
\subsection{Tailwind CSS -- Utility-First Styling}
\label{subsec:tailwind}

Traditionelles CSS erfordert separate Dateien und erfundene Klassennamen:

\begin{lstlisting}[style=base,numbers=none]
/* styles.css - waechst unkontrolliert */
.led-button { ... }
.led-button-active { ... }
.led-button-disabled { ... }
\end{lstlisting}

Tailwind bietet vordefinierte Utility-Klassen direkt im Markup:

\begin{lstlisting}[style=shell]
npm install -D tailwindcss postcss autoprefixer
npx tailwindcss init -p
\end{lstlisting}

\begin{lstlisting}[style=base,numbers=left,caption={Tailwind-Konfiguration},label={lst:tailwind-config}]
// tailwind.config.js
export default {
  content: ['./index.html', './src/**/*.{js,ts,jsx,tsx}'],
  theme: {
    extend: {
      colors: {
        'led-active': '#22c55e',
        'led-inactive': '#374151'
      }
    }
  }
};
\end{lstlisting}

\paragraph{Anwendung}

\begin{lstlisting}[style=base,numbers=left,caption={Tailwind-Klassen in React},label={lst:tailwind-react}]
// Statt CSS-Klassen zu erfinden: direkte Utilities
function LedButton({ isActive }: { isActive: boolean }) {
  return (
    <button className={`
      w-12 h-12 rounded-full
      transition-colors duration-150
      ${isActive
        ? 'bg-led-active shadow-lg shadow-green-500/50'
        : 'bg-led-inactive hover:bg-gray-600'}
    `}>
    </button>
  );
}
\end{lstlisting}

\paragraph{Grid für 100 LEDs}

\begin{lstlisting}[style=base,numbers=none]
<div className="grid grid-cols-10 gap-2 p-4">
  {leds.map((led, i) => <LedButton key={i} isActive={led} />)}
</div>
\end{lstlisting}

\begin{infobox}[Konsequenz]
Kein Wechsel zwischen Dateien, keine CSS-Spezifitätskonflikte. Die Bundle-Größe
bleibt klein, weil Tailwind nur genutzte Klassen einbindet.
\end{infobox}

% -----------------------------------------------------------------------------
\subsection{Zusammenspiel im Selection-Panel-Projekt}
\label{subsec:react-zusammenspiel}

\begin{figure}[H]
\centering
\begin{tikzpicture}[
  box/.style={draw, rectangle, rounded corners=2pt, minimum height=1.2cm,
              font=\ttfamily\footnotesize, fill=white, align=center},
  container/.style={draw, rectangle, rounded corners=3pt, minimum width=13cm,
                    minimum height=4cm, font=\sffamily\small},
  arrow/.style={-{Stealth[length=2mm]}, thick},
  label/.style={font=\sffamily\scriptsize, text=MutedText}
]
  % DEV Container
  \node[container, fill=ArduinoTeal!8] (dev) at (0,4) {};
  \node[font=\sffamily\bfseries, anchor=north west] at (-6.3,5.8) {Entwicklung (Mac)};

  \node[box] (vscode) at (-4,4.2) {VS Code\\+ Extensions};
  \node[box] (vite) at (0,4.2) {Vite Dev\\Server :5173};
  \node[box] (lint) at (4,4.2) {ESLint/\\Prettier};

  \draw[arrow] (vscode) -- (vite);
  \draw[arrow] (lint) -- (vite);

  \node[box] (bridge-dev) at (-4,2.5) {bridge.py\\:8081};
  \node[box, minimum width=2cm] (esp-dev) at (2,2.5) {ESP32-S3};

  \draw[arrow, dashed] (vite.south) -- node[right, label] {proxy /ws} (bridge-dev.north);
  \draw[arrow] (bridge-dev) -- node[above, label] {USB-CDC} (esp-dev);

  % PROD Container
  \node[container, fill=RaspberryRed!8, minimum height=2.5cm] (prod) at (0,-0.5) {};
  \node[font=\sffamily\bfseries, anchor=north west] at (-6.3,0.8) {Produktion (Raspberry Pi)};

  \node[box] (bridge-prod) at (-3,-0.5) {bridge.py\\:8081};
  \node[box, minimum width=2cm] (esp-prod) at (3,-0.5) {ESP32-S3};

  \node[label, anchor=east] at (0,-0.2) {liefert dist/};
  \node[label, anchor=east] at (0,-0.8) {WebSocket /ws};

  \draw[arrow] (bridge-prod) -- node[above, label] {USB-CDC} (esp-prod);
\end{tikzpicture}
\caption{Entwicklungs- und Produktionsumgebung}
\label{fig:dev-prod-setup}
\end{figure}

\paragraph{Deployment-Workflow}

\begin{lstlisting}[style=shell,caption={Frontend-Deployment},label={lst:frontend-deploy}]
# Auf dem Mac
cd apps/dashboard
npm run build          # Erzeugt dist/

# Auf den Pi kopieren
rsync -avh --delete \
  ./apps/dashboard/dist/ \
  pi@rover:/home/pi/embedded_projekt/apps/dashboard/dist/
\end{lstlisting}

% -----------------------------------------------------------------------------
\subsection{Siehe auch}
\label{subsec:react-siehe-auch}

Weiterführende Dokumentation:

\begin{itemize}
  \item \code{apps/dashboard/docs/DEVELOPER.md} -- Frontend-Entwicklung
  \item \code{apps/dashboard/docs/USER\_GUIDE.md} -- Benutzerhandbuch
  \item \code{apps/bridge/docs/PROTOCOL.md} -- WebSocket-Events
  \item \code{apps/bridge/docs/WORKFLOW.md} -- Pi-Deployment
\end{itemize}
